\documentclass{beamer}
\usetheme{Madrid}
\usecolortheme{beaver}

\usepackage[utf8]{inputenc}
\usepackage{amsmath}

\title{Red-Black Trees: Insertion and Rebalancing}
\subtitle{BLG 202E -- Data Structures}
\author{Deniz Sipahi -- 150230006}
\date{May 11, 2026}
\institute{Istanbul Technical University}

\begin{document}

% Slide 1 -- Title
\begin{frame}
\titlepage
\end{frame}

% Slide 2 -- Motivation
\begin{frame}{1. Motivation -- Why Balance Matters}
\begin{itemize}
    \item A \textbf{Binary Search Tree (BST)} supports search, insert, and delete in $O(h)$ time, where $h$ is the tree height.
    \item In the \textbf{worst case}, a BST degenerates into a linked list: $h = n$, so operations become $O(n)$.
    \item A \textbf{balanced BST} guarantees $h = O(\log n)$, making all operations efficient.
    \item \textbf{Red-Black Trees} are one of the most widely used self-balancing BSTs.
\end{itemize}

\begin{block}{Key Insight}
Balancing is not free --- it requires extra bookkeeping. Red-Black Trees achieve this through coloring constraints and rotations.
\end{block}
\end{frame}

% Slide 3 -- Red-Black Tree Properties
\begin{frame}{2. Red-Black Tree Properties}
Every Red-Black Tree satisfies these five invariants:
\begin{enumerate}
    \item Every node is either \textcolor{red}{\textbf{RED}} or \textbf{BLACK}.
    \item The root is always \textbf{BLACK}.
    \item Every leaf (NIL/null) is \textbf{BLACK}.
    \item If a node is \textcolor{red}{\textbf{RED}}, then both its children must be \textbf{BLACK}.
    \item For each node, all simple paths from the node to descendant leaves contain the same number of \textbf{BLACK} nodes.
\end{enumerate}
\end{frame}

% Slide 4 -- Node Structure
\begin{frame}{3. Node Structure}
Each node in a Red-Black Tree stores:
\begin{itemize}
    \item \textbf{Key}: the data value
    \item \textbf{Color}: RED or BLACK
    \item \textbf{Left}: pointer to the left child
    \item \textbf{Right}: pointer to the right child
    \item \textbf{Parent}: pointer to the parent node
\end{itemize}

\begin{block}{Representation}
\begin{tabular}{|c|c|c|c|c|}
\hline
\textbf{Key} & \textbf{Color} & \textbf{Left} & \textbf{Right} & \textbf{Parent} \\
\hline
\end{tabular}
\end{block}
\end{frame}

% Slide 5 -- Insertion Base Steps
\begin{frame}{4. Insertion -- Base Steps}
Inserting into a Red-Black Tree follows two phases:

\begin{enumerate}
    \item \textbf{Standard BST Insert}: Insert the new node as a leaf, colored \textcolor{red}{\textbf{RED}}.
    \item \textbf{Fix-up Phase}: Restore any violated RB-tree properties.
\end{enumerate}

\vspace{0.5em}
Possible violations after inserting a RED node:
\begin{itemize}
    \item \textbf{Property 2}: The new node might be the root (root must be BLACK).
    \item \textbf{Property 4}: The new node's parent might also be RED (no two consecutive RED nodes).
\end{itemize}
\end{frame}

% Slide 6 -- Fix-up Cases (ALL on one slide)
\begin{frame}{5. Fix-up Cases}
After insertion, we check the uncle of the new node:

\textbf{Case 1 -- Uncle is RED:}
Recolor the parent and uncle to BLACK, and the grandparent to RED. Move the pointer up to the grandparent and repeat.

\textbf{Case 2 -- Uncle is BLACK (zig-zag):}
The new node is on the opposite side of its parent relative to the grandparent. Perform a rotation on the parent to convert this into Case 3.

\textbf{Case 3 -- Uncle is BLACK (straight line):}
The new node is on the same side as its parent relative to the grandparent. Rotate on the grandparent and swap colors between the parent and grandparent.

\vspace{0.5em}
After fix-up, always recolor the root to BLACK.
\end{frame}

% Slide 7 -- Summary
\begin{frame}{6. Summary \& Key Takeaways}
\begin{itemize}
    \item Red-Black Trees maintain balance through \textbf{3 fix-up cases} after insertion.
    \item The fix-up process involves \textbf{recoloring} and at most \textbf{2 rotations}.
    \item All operations run in $O(\log n)$ time.
\end{itemize}

\begin{block}{Remember}
Think of the Red-Black Tree as a traffic system: RED means ``stop and check above,'' BLACK means ``all clear.'' Violations propagate upward like escalating a complaint to management.
\end{block}
\end{frame}

\end{document}
